فرض کنید یک جدول باید برای تابع 
\(f (x) = e^{x}\)
آماده شود، برای حالتی که 
\(x\)
در بازه [0, 1] است.
فرض کنید تعداد اعشاراتی که باید برای هر ورودی آمده باشد، 
\(d \geq 8\)
است، یعنی هر مقدار حداقل به هشت رقم اعشار نمایش داده می‌شود.
همچنین تفاوت بین مقادیر 
\(x\)
مجاور، یعنی اندازه گام، 
\(h\)
است. اندازه گام 
\(h\)
چقدر باید باشد تا تضمین شود که درون‌یابی خطی، خطای مطلق حداکثر 
\(10^{-6}\)
برای همه مقادیر
\(x\)
در بازه [0, 1] دارد؟

\begin{comment}
\end{comment}
پاسخ:

نقاط جدول را به صورت
\(x_j=jh\)
در نظر می‌گیریم. اگر
\(x_j\leq x\leq x_{j+1}\)
باشد، خطای درون‌یابی خطی برابر است با
\[
 f(x)-P_1(x)=\frac{f''(\xi)}{2}(x-x_j)(x-x_{j+1}),
 \qquad \xi\in[x_j,x_{j+1}].
\]
برای
\(f(x)=e^x\)
روی بازهٔ
\([0,1]\)
داریم
\(\max |f''|=e\).
همچنین بیشینهٔ
\(\left|(x-x_j)(x-x_{j+1})\right|\)
در نقطهٔ میانی زیر‌بازه برابر
\(h^2/4\)
است. پس خطای برش از کران زیر پیروی می‌کند:
\[
 |f(x)-P_1(x)|\leq\frac{eh^2}{8}.
\]

اگر مقادیر جدول دقیق فرض شوند، شرط
\(eh^2/8\leq10^{-6}\)
می‌دهد
\[
 h\leq\sqrt{\frac{8\times10^{-6}}{e}}
 \approx0.0017155278.
\]
برای یک شبکهٔ یکنواخت که دو سر بازه را نیز شامل کند باید
\(h=1/n\)
باشد؛ بنابراین کمترین تعداد زیر‌بازه
\(n=583\)
و بزرگ‌ترین گام مجاز در این حالت
\[
 h=\frac{1}{583}\approx0.0017152659
\]
است.

اگر عبارت «هشت رقم اعشار» به معنای گرد کردن مقادیر جدول تا هشت رقم اعشار باشد، خطای هر مقدار حداکثر
\(0.5\times10^{-8}\)
است. درون‌یابی خطی ترکیبی محدب از دو مقدار جدول است، پس همین مقدار کران خطای داده‌هاست و باید
\[
 \frac{eh^2}{8}+0.5\times10^{-8}\leq10^{-6}
\]
برقرار باشد. در این تفسیر کمترین تعداد زیر‌بازه
\(n=585\)
و بزرگ‌ترین گام شبکه
\[
 \boxed{h=\frac{1}{585}\approx0.0017094017}
\]
است.
